% Vietnamese translation of main.tex; references retain original publication titles.
% Internal research-in-progress draft; authors and affiliation are unconfirmed.
% Reference: D:/hoctap/hybird control system.pdf, supplied by the user.
% Elsevier CAS-based layout adapted to the measured production typography.
% No publisher identity, article DOI, acceptance dates or publication status
% from the sample is reproduced. This is not a journal-issued production proof.
\documentclass[fleqn]{cas-dc}
\usepackage[numbers,sort&compress]{natbib}
\usepackage{fontspec}
\setmainfont{Charis SIL}[
  Path=open_fonts/,
  UprightFont=CharisSIL-Regular.ttf,
  BoldFont=CharisSIL-Bold.ttf,
  ItalicFont=CharisSIL-Italic.ttf,
  BoldItalicFont=CharisSIL-BoldItalic.ttf
]
\geometry{paperwidth=210mm,paperheight=280mm,
  left=13.26mm,right=13.2mm,top=19.5mm,bottom=20mm,
  headsep=10pt,footskip=12pt,columnsep=18pt}
\hypersetup{colorlinks=true,citecolor=blue!55!black,
  urlcolor=blue!55!black,linkcolor=blue!55!black}
\urlstyle{same}
\renewcommand\refname{Tài liệu tham khảo}
\hypersetup{pdflang={vi-VN}}
\makeatletter
\renewcommand\normalsize{\@setfontsize\normalsize{8}{10.5}}
\renewcommand\small{\@setfontsize\small{7.2}{9}}
\renewcommand\footnotesize{\@setfontsize\footnotesize{6.4}{8.5}}
\renewcommand\scriptsize{\@setfontsize\scriptsize{6}{8}}
\normalsize
\renewcommand\sectionfont{\rmfamily\fontsize{8}{10.5}\bfseries\selectfont}
\renewcommand\ssectionfont{\rmfamily\fontsize{8}{10.5}\itshape\selectfont}
\renewcommand\printorcid{}
\AtBeginDocument{\renewcommand\bibfont{\fontsize{7.2}{9.5}\selectfont}}
\setlength{\parindent}{1em}
\setlength{\parskip}{0pt}
\newcommand{\DTAbstractFontAudit}{%
  \typeout{DTCHECK abstract size=\f@size pt; baseline=\the\baselineskip;
    family=\f@family; series=\f@series; shape=\f@shape}%
}
\def\ps@draftnotice{%
  \def\@oddhead{\hfil{\footnotesize\itshape Bản dịch tiếng Việt - Nghiên cứu đang triển khai}\hfil}%
  \let\@evenhead\@oddhead
  \def\@oddfoot{\hfil{\footnotesize\thepage}\hfil}%
  \let\@evenfoot\@oddfoot
}
\makeatother
% Adapt the public CAS front matter without editing the vendor class files.
% Fractions reproduce the supplied sample's narrow article-info panel.
\ExplSyntaxOn
\RenewDocumentEnvironment { Abstract } { o }
{
  \group_begin:
  \parindent0pt
  \box_if_empty:NTF \g_stm_key_box
    { \leftskip=.316\textwidth }
    {
      \dim_set:Nn \l_tmpa_dim { \box_ht:N\g_stm_key_box + \box_dp:N\g_stm_key_box }
      \leftskip=.316\textwidth
      \hspace*{-.316\textwidth}
      \noindent\hbox_to_wd:nn { 0pt } { \box\g_stm_key_box\hss }
      \skip_vertical:n { -\l_tmpa_dim }
    }
  \noindent{\fontsize{7}{9}\selectfont TÓM~TẮT}\par
  \skip_vertical:n { -4pt }
  \noindent\rule{.684\textwidth}{.2pt}\par
  \fontsize{7.2}{9}\selectfont
  \DTAbstractFontAudit
  \ignorespaces
}
{ \par\group_end: }
\RenewDocumentEnvironment { keywords } { O{Từ~khóa} }
{
  \tex_global:D\tex_setbox:D\g_stm_key_box=\vtop\bgroup
  \hsize=.256\textwidth
  \cs_set:Npn\sep{\par}
  \parindent0pt
  {\fontsize{7}{9}\selectfont THÔNG~TIN~BÀI~VIẾT}\par
  \skip_vertical:n{-3pt}
  \rule{.256\textwidth}{.2pt}\par\footnotesize
  \noindent\textit{#1}:\par
}
{ \egroup }
\ExplSyntaxOff
\begin{document}
\shorttitle{Bản sao số dự báo có nhận thức về mạng}
\shortauthors{}
\title[mode=title,size={\fontsize{13.5}{17.2}\selectfont},after=15pt]
{Khung bản sao số dự báo có nhận thức về mạng cho hệ thống sản xuất không gian mạng - vật lý có khả năng chống chịu và tự tái cấu hình}
% Blank, reserved fields: user has not yet supplied an author list or affiliation.
\author{\mbox{}}
\begin{abstract}
Lỗi cảm biến và thông tin trạng thái bị trễ có thể khiến bộ điều khiển sản xuất và biểu diễn số của nó không thống nhất về luồng vật liệu và khả năng sẵn sàng của thiết bị. Nghiên cứu này đề xuất một khung bản sao số dự báo có nhận thức về mạng để tái cấu hình ở cấp phần mềm cho trạm phân loại và gắp - thả sử dụng băng tải. Mô hình MuJoCo gồm hai băng tải, cơ cấu chuyển phôi khí nén, robot tọa độ Descartes và cơ cấu gắp chân không tạo thành nền tảng mô phỏng ban đầu. Kiến trúc đề xuất tách thông tin quan sát có dấu thời gian từ cảm biến khỏi trạng thái chuẩn của bộ mô phỏng, đồng thời kết hợp ước lượng trạng thái lai, phát hiện bất thường dựa trên phần dư và mô phỏng các kịch bản dự báo chịu ràng buộc về hạn chót tính toán. Độ trễ truyền thông, độ biến thiên trễ, mất gói tin và độ mới của dữ liệu được sử dụng làm đầu vào để xác định độ tin cậy của trạng thái và các quyết định giám sát. Các đáp ứng ứng viên gồm điều chỉnh thời điểm vận hành băng tải, hạn chế đưa phôi vào, sử dụng thông tin quan sát hợp lệ còn lại và tạm dừng hoạt động khi độ bất định vượt quá giới hạn phụ thuộc nhiệm vụ. Việc tái cấu hình chỉ giới hạn ở các chính sách được cơ cấu chấp hành đã lắp đặt và các liên động điều khiển cục bộ hỗ trợ. Kế hoạch đánh giá so sánh điều khiển tuần tự, giám sát có xét lỗi và giám sát dự báo có nhận thức về mạng với cùng lịch tác động lỗi và điều kiện vận hành. Các tiêu chí gồm sai số ước lượng trạng thái, độ trễ phát hiện lỗi, tần suất báo động giả, khả năng hoàn thành nhiệm vụ, thời gian phục hồi và khả năng đáp ứng hạn chót dự báo. Nguyên mẫu hiện tại cung cấp môi trường mô phỏng vật lý; bộ ước lượng, bộ giám sát có nhận thức về mạng và việc kiểm chứng định lượng vẫn đang được phát triển. Mục tiêu nghiên cứu là xác định khi nào giám sát dự báo có thể duy trì tính liên tục của sản xuất trước các nhiễu cảm biến và truyền thông mà không thay đổi cấu hình cơ khí.
\end{abstract}
\begin{keywords}
Hệ thống sản xuất không gian mạng - vật lý \sep
Bản sao số \sep
Chẩn đoán lỗi \sep
Điều khiển có nhận thức về mạng \sep
Tái cấu hình phần mềm
\end{keywords}
\maketitle
\pagestyle{draftnotice}
\thispagestyle{draftnotice}
\normalsize
\section{Giới thiệu}
% Bản dịch tương ứng với C101--C109 và các nguồn của introduction.tex.
Các hệ thống sản xuất không gian mạng - vật lý kết nối thiết bị vận chuyển vật liệu, cảm biến và phần mềm điều khiển để phối hợp hoạt động vật lý với các quyết định tính toán. Trong trạm sản xuất được xem xét ở đây, khả năng duy trì sản xuất phụ thuộc vào việc xác định phôi đã đến vị trí chuyển tiếp hay chưa, chuyển động khí nén đã hoàn tất hay chưa và robot có thể thực hiện thao tác tiếp theo hay không. Một biểu diễn số trở nên hữu ích cho mục đích này khi các ước lượng trạng thái và dự báo của nó có thể hỗ trợ những quyết định đó. Kritzinger và cộng sự phân loại các biểu diễn trong sản xuất thành mô hình số, bóng số và bản sao số theo mức độ tích hợp~\cite{Kritzinger2018}. Sự phân biệt này định hướng phạm vi của nghiên cứu: mô hình vật lý có thể thực thi tạo thành nền tảng mô phỏng ban đầu, còn xử lý thông tin quan sát, đồng bộ hóa, dự báo và tác động giám sát được đề xuất dưới dạng các lớp bổ sung.

Ứng dụng mục tiêu là một trạm phân loại và gắp - thả gồm hai băng tải, một xi lanh đẩy khí nén, một robot tọa độ Descartes và một cơ cấu gắp chân không. Cấu hình cảm biến dự kiến bao gồm phát hiện phôi bằng cảm biến quang và phản hồi vị trí cuối hành trình của chuyển động khí nén. Bộ điều khiển tuần tự lần lượt chuyển qua các trạng thái cấp phôi, chuyển tiếp, theo dõi, gắp, đặt và trở về. Trong ứng dụng này, việc thiếu một sự kiện phát hiện có thể xuất phát từ nhiều nguyên nhân: phôi chưa đến, cảm biến bị hỏng hoặc thông tin quan sát bị trễ hay mất trong quá trình truyền. Vì vậy, đáp ứng giám sát cần xét đồng thời hành vi của thiết bị và tính hợp lệ của thông tin dùng để suy luận hành vi đó. Tiếp tục chu trình dựa trên một ước lượng trạng thái thiếu căn cứ và dừng mỗi khi thiếu thông tin quan sát là hai lựa chọn vận hành khác nhau; hệ quả của chúng cần được đánh giá dưới cùng các nhiễu tác động.

Ảnh hưởng của truyền thông là một phần của bài toán điều khiển khi thông tin trạng thái được truyền qua mạng. Hespanha và cộng sự tổng quan các phương pháp ước lượng và điều khiển dưới những hạn chế như lấy mẫu, độ trễ biến thiên và mất gói tin~\cite{Hespanha2007}. Yates và cộng sự xem xét tuổi thông tin như một đại lượng đo tính kịp thời của các bản cập nhật trạng thái có dấu thời gian~\cite{Yates2021}. Những cách tiếp cận này là cơ sở để coi tuổi và lịch sử truyền nhận của thông tin quan sát là đầu vào của bản sao số được đề xuất. Với trạm hiện tại, ước lượng vị trí phôi suy ra từ một sự kiện cảm biến quang đã cũ cần có mức độ tin cậy khác với ước lượng được hiệu chỉnh bằng một sự kiện mới. Do đó, khi triển khai, hệ thống sẽ phân biệt dấu thời gian tại nguồn, thời điểm tiếp nhận, số thứ tự và các miền đồng hồ. Độ bất định của việc đồng bộ đồng hồ sẽ được tính đến khi diễn giải tuổi của thông tin quan sát, thay vì giả định dấu thời gian từ các thiết bị khác nhau có chung một cơ sở thời gian chính xác.

Khả năng thích ứng với lỗi cảm biến cũng đã được nghiên cứu trong sản xuất tự động. El-Mahdy và cộng sự mô tả một phương pháp chịu lỗi chủ động, sử dụng các mạng bộ nhớ dài ngắn hạn riêng biệt cho chẩn đoán và tái cấu hình điều khiển dựa trên mô hình nghịch đảo; phương pháp được đánh giá bằng mô phỏng vận chuyển vật liệu trên Factory I/O--MATLAB~\cite{ElMahdy2024}. Nghiên cứu đó cung cấp một ví dụ phù hợp về việc liên kết chẩn đoán với điều chỉnh tác động điều khiển. Trong dự án hiện tại, thiết kế ban đầu sẽ sử dụng các phần dư có thể diễn giải, kiểm tra tính nhất quán của sự kiện và các giới hạn bất định tường minh. Lựa chọn này cho phép liên hệ một sự kiện phôi bị bỏ sót, một chuyển trạng thái xi lanh quá hạn hoặc dữ liệu đã cũ với một giả định vận hành cụ thể, trước khi xem xét các mô hình học máy phức tạp hơn.

Điều khiển dự báo bằng bản sao số cung cấp một cơ sở bổ sung để đánh giá các hành động trong tương lai. Chen và cộng sự kết hợp dự báo nhiều bước và ước lượng bất định dựa trên phân vị với điều khiển dự báo mô hình bền vững, trong đó có một nghiên cứu tình huống về sản xuất bồi đắp~\cite{Chen2026}. Công trình đó gợi mở việc xem xét đồng thời độ bất định của dự báo và thời gian tính toán. Trong khi đó, bài toán gắp - thả ở đây tập trung vào thời điểm xảy ra các sự kiện sản xuất rời rạc và việc lựa chọn trong một tập hữu hạn các chính sách giám sát. Câu hỏi thiết kế là liệu dự báo còn hữu ích hay không khi thông tin quan sát dùng để khởi tạo dự báo có thể bị lỗi hoặc đã cũ, đồng thời đáp ứng được chọn phải được trả về trước quyết định gắp - thả tiếp theo.

Từ đó, câu hỏi nghiên cứu được đặt ra như sau: trong những điều kiện cảm biến và truyền thông nào, một bản sao số dự báo có nhận thức về mạng có thể duy trì khả năng hoàn thành nhiệm vụ bằng cách lựa chọn một chính sách phần mềm được phép trước hạn chót ra quyết định? Trả lời câu hỏi này đòi hỏi kết hợp độ tin cậy của trạng thái, các giả thuyết lỗi, chi phí dự báo và các tác động điều khiển sẵn có. Chỉ một mô phỏng chi tiết về hình học là chưa đủ để trả lời. Bài toán tích hợp cụ thể mà dự án giải quyết là truyền tải độ bất định gắn với thông tin quan sát nhận được xuyên suốt quá trình chẩn đoán và đánh giá kịch bản trong khoảng dự báo ngắn, để đi đến quyết định giám sát. Nghiên cứu cũng sẽ xem xét khi nào các thông tin quan sát còn lại không đủ để tiếp tục thao tác và việc tạm dừng có kiểm soát là đáp ứng phù hợp.

Để khảo sát câu hỏi này, khung đề xuất được xây dựng trên mô hình trạm MuJoCo hiện có. MuJoCo cung cấp các chức năng động lực học đa vật thể và tiếp xúc phục vụ điều khiển dựa trên mô hình~\cite{Todorov2012}. Trong kiến trúc thực nghiệm dự kiến, một phiên bản bộ mô phỏng đóng vai trò đối tượng sản xuất, còn một phiên bản dự báo độc lập được khởi tạo từ các ước lượng suy ra qua cảm biến. Trạng thái chuẩn của bộ mô phỏng chỉ được dành cho đánh giá. Bộ ước lượng lai sẽ kết hợp các ước lượng chuyển động liên tục với trạng thái nhiệm vụ rời rạc; lớp giám sát lỗi sẽ đánh giá phần dư cùng với thứ tự sự kiện, thời gian chờ và độ mới của dữ liệu. Các lượt mô phỏng dự báo sẽ khảo sát những nhiễu có thể xảy ra và các đáp ứng được phép trong một ngân sách tính toán cố định. Bộ giám sát chỉ phát hành đáp ứng khi các giả định về trạng thái và khoảng thời gian hiệu lực của lệnh vẫn còn phù hợp. Bộ điều khiển trình tự và các liên động cục bộ vẫn giữ quyền quyết định việc thực thi tại cơ cấu chấp hành.

Ba đóng góp được đề xuất gồm:
\begin{enumerate}
  \item Một kiến trúc liên kết thông tin quan sát có dấu thời gian, ước lượng trạng thái lai, giả thuyết lỗi và độ tin cậy của thông tin quan sát với các quyết định giám sát cho trạm vận chuyển và gắp - thả sử dụng băng tải.
  \item Một cơ chế dự báo và lựa chọn nhằm đánh giá các chính sách vận hành được phép dưới tác động của nhiễu cảm biến và truyền thông, đồng thời tuân thủ hạn chót tính toán và duy trì một đáp ứng tạm dừng đã được xác định.
  \item Một quy trình đánh giá so sánh điều khiển tuần tự, giám sát có xét lỗi và giám sát dự báo có nhận thức về mạng với cùng điều kiện ban đầu và lịch tác động nhiễu; các chỉ tiêu được tách riêng cho chất lượng ước lượng, hiệu năng phát hiện, khả năng hoàn thành sản xuất, phục hồi và tính kịp thời của tính toán.
\end{enumerate}

Trong nghiên cứu này, tự tái cấu hình là thay đổi các chính sách phần mềm được thiết bị hiện có hỗ trợ: điều chỉnh thời điểm vận hành băng tải, hạn chế đưa phôi mới vào, lựa chọn thông tin quan sát hợp lệ hoặc tạm dừng chu trình. Điều chỉnh tốc độ liên tục chỉ là một đáp ứng được phép khi bộ truyền động băng tải đã lắp đặt hỗ trợ chức năng đó. Nguyên mẫu hiện tại cung cấp mô phỏng vật lý và một trình tự chạy thử; bộ ước lượng, mô hình truyền thông, bộ giám sát dự báo và việc kiểm chứng tích hợp vẫn đang được phát triển. Việc hiệu chuẩn theo trạm vật lý cũng chưa hoàn tất. Vì vậy, bản thảo nghiên cứu đang triển khai này trình bày cách đặt bài toán và các đóng góp phương pháp dự kiến, không khẳng định đã đo được mức cải thiện về thông lượng, độ chính xác phát hiện lỗi hoặc hiệu năng phục hồi.

\balance
% Numeric Elsevier-style references. Metadata: Crossref; access scope:
% evidence/crossref_metadata.json and evidence/introduction_sources.md.
\begin{thebibliography}{6}
\bibitem{Kritzinger2018}
W.~Kritzinger, M.~Karner, G.~Traar, J.~Henjes, W.~Sihn,
Digital twin in manufacturing: A categorical literature review and classification,
\emph{IFAC-PapersOnLine} 51~(11) (2018) 1016--1022,
\url{https://doi.org/10.1016/j.ifacol.2018.08.474}.

\bibitem{Hespanha2007}
J.~P.~Hespanha, P.~Naghshtabrizi, Y.~Xu,
A survey of recent results in networked control systems,
\emph{Proceedings of the IEEE} 95~(1) (2007) 138--162,
\url{https://doi.org/10.1109/JPROC.2006.887288}.

\bibitem{Yates2021}
R.~D.~Yates, Y.~Sun, D.~R.~Brown~III, S.~K.~Kaul, E.~Modiano, S.~Ulukus,
Age of information: An introduction and survey,
\emph{IEEE Journal on Selected Areas in Communications} 39~(5) (2021) 1183--1210,
\url{https://doi.org/10.1109/JSAC.2021.3065072}.

\bibitem{ElMahdy2024}
M.~H.~El-Mahdy, M.~I.~Awad, S.~A.~Maged,
Deep-learning-based design of active fault-tolerant control for automated manufacturing systems subjected to faulty sensors,
\emph{Transactions of the Institute of Measurement and Control} 46~(12) (2024) 2289--2299,
\url{https://doi.org/10.1177/01423312241229493}.

\bibitem{Chen2026}
Y.-P.~Chen, Y.-K.~Tsai, V.~Karkaria, W.~Chen,
Uncertainty-aware digital twins: Robust model predictive control using time-series deep quantile learning,
\emph{Journal of Mechanical Design} 148~(2) (2026) 021702,
\url{https://doi.org/10.1115/1.4069104}.

\bibitem{Todorov2012}
E.~Todorov, T.~Erez, Y.~Tassa,
MuJoCo: A physics engine for model-based control,
in: \emph{2012 IEEE/RSJ International Conference on Intelligent Robots and Systems},
2012, pp.~5026--5033,
\url{https://doi.org/10.1109/IROS.2012.6386109}.
\end{thebibliography}

\end{document}

